\chapter{Durchführung} \label{ch:Durchfuehrung}

\section{Messverfahren} \label{sec:Messverfahren}

Zu Beginn des Versuchs wird mithilfe der Messsoftware eine Offsetmessung durchgeführt, um den Gleichgewichtswinkel beider Pendel zu bestimmen. Dabei müssen die Pendel vollständig in Ruhe sein und die Kopplungsfeder darf noch nicht montiert sein. Der gemessene Offset wird anschließend automatisch von allen folgenden Messdaten abgezogen.

Im ersten Teil des Experiments wird die Eigenfrequenz der ungekoppelten Pendel bestimmt. Dazu werden beide Pendel leicht ausgelenkt und frei schwingen gelassen. Die Software zeichnet die Winkelauslenkung über die Zeit auf und berechnet mittels FFT die Schwingungsfrequenz. Da beide Pendel baugleich sind, genügt eine gemeinsame Messung.

Anschließend wird die Kopplungsfeder montiert. Für drei verschiedene Aufhängepunkte der Feder werden jeweils die symmetrische und die antisymmetrische Eigenschwingung angeregt. Bei der symmetrischen Anregung werden beide Pendel gleich weit und gleichsinnig ausgelenkt, bei der antisymmetrischen gegensinnig. Die Frequenzen der jeweiligen Normalschwingungen werden erneut über die FFT bestimmt.

Für die Untersuchung der Schwebung wird eines der Pendel ausgelenkt, während das andere in Ruhe gehalten wird. Nach dem gleichzeitigen Loslassen entsteht eine Überlagerung beider Eigenmoden, die sich in einer periodischen Energieübertragung zwischen den Pendeln äußert. Neben den beiden Eigenfrequenzen werden hier zusätzlich die Schwingungs- und Schwebungsperioden direkt aus dem Zeitverlauf mittels Cursormessung bestimmt.

Im letzten Teil des Versuchs wird ein analoges System aus zwei induktiv gekoppelten elektrischen Schwingkreisen untersucht. Durch Variation des Abstands der Spulen wird die Kopplungsstärke verändert. Die Spannungsverläufe beider Kreise werden mit einem Oszilloskop beobachtet, wobei insbesondere das Auftreten von Schwebungen qualitativ analysiert wird.

\img[0.25]{memes/bLatex}{Ja Dimitrij, ich kenne dich...}